\documentclass[a4paper,14pt]{extarticle}

\usepackage[left=30mm,right=10mm,top=20mm,bottom=20mm]{geometry}
\usepackage{iftex}
\ifPDFTeX
  \usepackage[T2A]{fontenc}
  \usepackage[utf8]{inputenc}
  \usepackage[english,russian]{babel}
  % Keep default CM fonts for stable Cyrillic support in pdfLaTeX.
\else
  \usepackage{fontspec}
  \usepackage{polyglossia}
  \setdefaultlanguage{russian}
  \setotherlanguage{english}
  % Strict mode for methodology: if Times New Roman TTF files are uploaded to
  % fonts/, use them; otherwise fallback to CMU fonts for guaranteed compile.
  \IfFileExists{fonts/times.ttf}{%
    \setmainfont{times.ttf}[
      Path=fonts/,
      BoldFont=timesbd.ttf,
      ItalicFont=timesi.ttf,
      BoldItalicFont=timesbi.ttf,
      Script=Cyrillic,
      Ligatures=TeX
    ]%
    \newfontfamily\cyrillicfont{times.ttf}[
      Path=fonts/,
      BoldFont=timesbd.ttf,
      ItalicFont=timesi.ttf,
      BoldItalicFont=timesbi.ttf,
      Script=Cyrillic
    ]%
    \setsansfont{arial.ttf}[
      Path=fonts/,
      BoldFont=arialbd.ttf,
      ItalicFont=ariali.ttf,
      BoldItalicFont=arialbi.ttf,
      Script=Cyrillic,
      Ligatures=TeX
    ]%
    \newfontfamily\cyrillicfontsf{arial.ttf}[
      Path=fonts/,
      BoldFont=arialbd.ttf,
      ItalicFont=ariali.ttf,
      BoldItalicFont=arialbi.ttf,
      Script=Cyrillic
    ]%
    \setmonofont{cour.ttf}[
      Path=fonts/,
      BoldFont=courbd.ttf,
      ItalicFont=couri.ttf,
      BoldItalicFont=courbi.ttf,
      Script=Cyrillic,
      Ligatures=TeX
    ]%
    \newfontfamily\cyrillicfonttt{cour.ttf}[
      Path=fonts/,
      BoldFont=courbd.ttf,
      ItalicFont=couri.ttf,
      BoldItalicFont=courbi.ttf,
      Script=Cyrillic
    ]%
  }{%
    \setmainfont{CMU Serif}[Script=Cyrillic,Ligatures=TeX]
    \newfontfamily\cyrillicfont{CMU Serif}[Script=Cyrillic]
    \setsansfont{CMU Sans Serif}[Script=Cyrillic,Ligatures=TeX]
    \newfontfamily\cyrillicfontsf{CMU Sans Serif}[Script=Cyrillic]
    \setmonofont{CMU Typewriter Text}[Script=Cyrillic,Ligatures=TeX]
    \newfontfamily\cyrillicfonttt{CMU Typewriter Text}[Script=Cyrillic]
  }
\fi
\AtBeginDocument{\selectlanguage{russian}}

\usepackage{setspace}
\onehalfspacing
\setlength{\parindent}{1.25cm}
\setlength{\parskip}{0pt}
% Reduce noisy box warnings to avoid huge logs/timeouts on Overleaf free plan.
\setlength{\emergencystretch}{3em}
\hfuzz=30pt
\hbadness=10000
\vbadness=10000
\sloppy

\usepackage{titlesec}
\titleformat{\section}{\bfseries\centering}{\thesection}{1em}{}
\titleformat{name=\section,numberless}{\bfseries\centering}{}{0pt}{}
\titleformat{\subsection}{\bfseries\centering}{\thesubsection}{1em}{}
\titlespacing*{\section}{0pt}{0pt}{2\baselineskip}
\titlespacing*{\subsection}{0pt}{2\baselineskip}{2\baselineskip}

\usepackage{fancyhdr}
\fancypagestyle{main}{
  \fancyhf{}
  \fancyfoot[C]{\thepage}
  \renewcommand{\headrulewidth}{0pt}
  \renewcommand{\footrulewidth}{0pt}
}

\usepackage{graphicx}
\usepackage{booktabs}
\usepackage{array}
\usepackage{longtable}
\usepackage{caption}
\usepackage{amsmath}
\usepackage{amssymb}
\usepackage{chngcntr}
\counterwithin{table}{section}
\counterwithin{figure}{section}
\numberwithin{equation}{section}

\captionsetup[figure]{justification=centering,labelsep=period}
\captionsetup[table]{justification=raggedright,singlelinecheck=false,labelsep=period}

\usepackage[round,authoryear]{natbib}
\setcitestyle{authoryear,aysep={, },yysep={;},notesep={, }}
\let\cite\citep
\usepackage{url}

\addto\captionsrussian{%
  \renewcommand{\contentsname}{СОДЕРЖАНИЕ}%
  \renewcommand{\figurename}{Рисунок}%
  \renewcommand{\tablename}{Таблица}%
}
\addto\captionsenglish{%
  \renewcommand{\contentsname}{СОДЕРЖАНИЕ}%
}
\renewcommand{\refname}{СПИСОК ИСПОЛЬЗОВАННОЙ ЛИТЕРАТУРЫ}

\newcommand{\appendixsection}[2]{%
  \clearpage
  \thispagestyle{empty}
  \noindent\hfill\textbf{Приложение #1}

  \vspace{1.5\baselineskip}
  \begin{center}
    \textbf{#2}
  \end{center}
  \vspace{\baselineskip}
}

\begin{document}

\begin{titlepage}
\thispagestyle{empty}
\begin{center}
Министерство просвещения Республики Казахстан \\
Колледж ТОО <<Astana IT University>> \\
Цикловая комиссия <<Специальных дисциплин>>

\vspace{4\baselineskip}
\textbf{КУРСОВОЙ ПРОЕКТ}

\vspace{2\baselineskip}
по дисциплине: \textbf{Разработка мобильных приложений}

\vspace{1.2\baselineskip}
на тему: \textbf{Разработка CRM-приложения Gamified Task CRM}

\vfill
\begin{flushright}
\begin{minipage}{0.46\textwidth}
\raggedright
Выполнил: Тойшыбаев Жаслан\\[0.6\baselineskip]
студент 3 курса, группа ПО2303\\[0.6\baselineskip]
Руководитель: Арсен Тимурович\\[0.6\baselineskip]
должность: Преподаватель по специальным дисциплинам
\end{minipage}
\end{flushright}

\vfill
Астана, 2026
\end{center}
\end{titlepage}

\clearpage
\thispagestyle{empty}
\tableofcontents

\clearpage
\setcounter{page}{3}
\pagestyle{main}

\section*{ВВЕДЕНИЕ}
\addcontentsline{toc}{section}{ВВЕДЕНИЕ}

В условиях устойчивой цифровизации экономики и социальной сферы все больше процессов управления взаимодействием с клиентами переводится в гибридный и полностью цифровой формат. Для учебных проектов по направлению разработки программного обеспечения это означает смещение фокуса от простых демонстрационных приложений к системам, которые отражают реальный жизненный цикл продукта: учет требований, архитектурное проектирование, разработка пользовательского интерфейса, интеграция с серверной частью, контроль качества и подготовка к развертыванию. В этой логике система \textit{Gamified Task CRM (Flutter Web)} рассматривается как прикладной программный продукт, объединяющий функции управления задачами, базовые CRM-процессы, модуль отчетности и вспомогательные интеллектуальные сервисы \cite{repoReadme,docsIndex}.

Актуальность темы курсового проекта определяется сразу несколькими факторами. Во-первых, бизнес-пользователи ожидают от CRM-решений не только хранения карточек клиентов, но и связности всех операционных сущностей: задач, этапов работы, аналитики и персональных метрик. Во-вторых, активно растет спрос на кроссплатформенные технологии, позволяющие быстро выпускать единый продукт для web-окружения без отдельной реализации на нескольких стековых решениях \cite{flutterDocs,dartDocs}. В-третьих, в образовательном процессе становится критичным умение строить приложение на современных практиках: декларативный UI, управление состоянием, строгая маршрутизация, безопасная работа с конфигурацией и воспроизводимый тестовый контур \cite{riverpodDocs,goRouterDocs,owaspM10}. Таким образом, выбранная тема позволяет связать теоретические положения дисциплины с практически значимой реализацией.

Дополнительный аргумент в пользу выбранной темы состоит в том, что разрабатываемая система находится на стыке нескольких предметных областей: проектирование интерфейсов, управление данными, прикладная аналитика, инженерия качества и основы информационной безопасности. Именно такие комплексные проекты позволяют продемонстрировать уровень профессиональной подготовки студента, поскольку требуют не фрагментарных знаний, а целостного подхода к разработке программного обеспечения \cite{sommerville,cleanArchitecture}.

\textbf{Целью} курсового проекта является разработка и обоснование архитектурно-целостного web-приложения \textit{Gamified Task CRM} для управления задачами и клиентскими взаимодействиями с использованием современного Flutter-стека и backend-first интеграции с Supabase.

Для достижения поставленной цели в работе решаются следующие задачи:
\begin{enumerate}
  \item Проанализировать теоретические и технологические основы построения CRM-ориентированных систем управления задачами, включая требования к архитектуре, данным и пользовательскому опыту.
  \item Спроектировать логическую структуру программного решения: модули интерфейса, маршрутизацию, модель состояния, состав ключевых сущностей и взаимодействие с backend-сервисами.
  \item Реализовать практическую часть проекта на Flutter Web с поддержкой аутентификации, CRUD-операций, связей \textit{задача-клиент}, базовой аналитики и конфигурационных режимов выполнения.
  \item Проверить корректность работы решения с использованием встроенных сценариев качества (analyze, unit/widget/integration тесты, smoke/E2E-потоки) и сформулировать итоговые выводы о достигнутых результатах.
\end{enumerate}

Объектом исследования в курсовом проекте выступают процессы цифрового управления задачами и клиентскими взаимодействиями в образовательном и малом бизнес-контексте. Предмет исследования --- методы и средства разработки кроссплатформенного CRM-приложения на базе Flutter, Riverpod, GoRouter и Supabase с учетом эксплуатационных требований к надежности, расширяемости и удобству использования \cite{supabaseDocs,postgresDocs}.

Практическая значимость проекта заключается в том, что реализованное приложение может применяться как учебно-прикладной стенд для демонстрации полного цикла разработки: от постановки требований до тестирования и подготовки к деплою. Кроме того, сформированная архитектура допускает масштабирование в сторону более сложных бизнес-сценариев: сегментации клиентов, автоматизации воронки продаж, расширенной отчетности, уведомлений и интеграции внешних каналов взаимодействия \cite{crmBook,projectManagementBook}.

Методологическую основу работы составляют: системный анализ требований, модульный принцип проектирования, сравнительный анализ технологических подходов, итеративная реализация и верификация качества через тестовые сценарии. При подготовке курсового проекта использовались официальные технические руководства по Flutter, Dart, Riverpod, GoRouter и Supabase, а также учебные и научно-практические источники по проектированию информационных систем и программной архитектуре \cite{flutterDocs,dartDocs,riverpodDocs,goRouterDocs,supabaseDocs,isDesign2025}.

Ключевым организационно-техническим принципом проекта является разделение рабочих режимов выполнения. В базовом сценарии используется backend-first подход с реальной авторизацией и операциями чтения/записи в Supabase. Для учебной демонстрации также предусмотрен локальный demo-режим с подстановкой тестовых данных, что позволяет воспроизводить функционал даже при отсутствии внешнего окружения. Такое решение повышает устойчивость проекта в условиях лабораторной и защитной демонстрации, но при этом требует явного разграничения демо и production-потоков для минимизации риска смешения данных \cite{repoReadme,owaspM4}.

Отдельное внимание в работе уделяется качеству программного продукта. Проверка корректности не ограничивается ручным тестированием пользовательского интерфейса. В проект включены команды статического анализа, модульные и виджет-тесты, интеграционные smoke-сценарии и web E2E-проверка основного бизнес-потока \textit{Login $\rightarrow$ Create Client $\rightarrow$ Create Task $\rightarrow$ Verify bindings}. Такой набор позволяет обнаруживать ошибки на разных уровнях системы и снижает вероятность регрессий при дальнейшем развитии кода \cite{repoReadme,testingBook}.

С точки зрения структуры, курсовой проект состоит из введения, двух разделов основной части, заключения, списка использованной литературы и приложений. В первом разделе рассматриваются теоретические основы разработки CRM-ориентированных приложений и обосновывается выбор технологического стека. Во втором разделе представлена практическая реализация \textit{Gamified Task CRM}: архитектура, модель данных, ключевые пользовательские сценарии и результаты верификации. В заключении подведены итоги, дана оценка степени достижения цели и обозначены направления дальнейшего развития системы.

В рамках ограничений курсового проекта акцент сделан на базовом эксплуатационном контуре: корректная работа аутентификации, управление задачами, учет клиентов, формирование отчетных представлений и воспроизводимость CI-проверок. Расширенные возможности (глубокая аналитика, многоканальные уведомления, продвинутая сегментация) рассматриваются как логические точки роста и не противоречат текущей архитектуре.

Таким образом, выбранная тема и реализуемый программный продукт соответствуют требованиям учебной дисциплины, обладают практической применимостью и позволяют в полной мере продемонстрировать навыки системного проектирования и разработки программного обеспечения на современном технологическом стеке.

\clearpage
\section{ТЕОРЕТИЧЕСКИЕ ОСНОВЫ РАЗРАБОТКИ CRM-ОРИЕНТИРОВАННОГО ПРИЛОЖЕНИЯ}

\subsection{Предметная область и функциональные требования к CRM-системам}

Современная CRM-система в прикладном понимании представляет собой информационную платформу, обеспечивающую централизованный учет клиентов, историю взаимодействия, планирование активностей и контроль результативности операций. Если рассматривать учебный проект, то минимально достаточная модель должна включать сущности клиента, задачи, статусы выполнения, источники лидов и пользовательскую аналитику. Без этих компонентов невозможно обеспечить трассируемость бизнес-действий и сформировать корректную отчетную базу \cite{crmBook,projectManagementBook}.

Классическая задача CRM в цифровой среде состоит в том, чтобы перевести разрозненные процессы коммуникации в формализованные сценарии: от первичного контакта до конверсии и постпродажного сопровождения. На уровне данных это приводит к необходимости поддерживать консистентные связи между сущностями. Например, задача должна быть либо автономной, либо привязанной к конкретному клиенту, а изменение статуса клиента должно корректно отражаться в аналитических представлениях. В противном случае отчетность становится недостоверной, а управленческие решения теряют обоснованность \cite{dataDesign2022,isDesign2023}.

Для курсового проекта важным является ограничение функционального объема. С одной стороны, работа должна быть достаточно содержательной, чтобы продемонстрировать инженерные компетенции. С другой стороны, она не должна превращаться в попытку реализации корпоративной ERP-системы. Поэтому в рамках \textit{Gamified Task CRM} функциональные требования сгруппированы вокруг четырех опорных блоков:
\begin{enumerate}
  \item Пользовательский контур: авторизация, переходы между основными разделами приложения, обратная связь по ошибкам.
  \item Операционный контур: управление задачами и клиентами, поддержка CRUD-операций, назначение приоритетов, учет сроков.
  \item Аналитический контур: представление агрегированной информации по активности и прогрессу пользователя.
  \item Эксплуатационный контур: режимы выполнения (demo и backend-first), воспроизводимые проверки качества и деплой web-версии.
\end{enumerate}

Такое деление соответствует принципу функциональной декомпозиции, когда каждую подсистему можно развивать независимо при сохранении общих интерфейсных контрактов. Применение этого подхода в учебных проектах повышает управляемость кода и упрощает тестирование, поскольку сценарии проверки привязываются к четко выделенным зонам ответственности \cite{cleanArchitecture,patternsBook}.

Отдельно необходимо учитывать нефункциональные требования, которые часто игнорируются на ранних этапах обучения, но критически важны при защите проекта. К ключевым нефункциональным требованиям относятся:
\begin{itemize}
  \item надежность операций записи и чтения данных;
  \item предсказуемость пользовательских сценариев при ошибках сети и backend;
  \item масштабируемость структуры кода при добавлении новых модулей;
  \item безопасность хранения и передачи конфигурационных параметров;
  \item прозрачность контроля качества через автоматизированные проверки.
\end{itemize}

С методической точки зрения указанные требования создают мост между теорией системного анализа и практикой программной инженерии. Именно их учет в курсовом проекте позволяет перейти от формального описания интерфейсов к аргументированному выбору архитектуры.

\subsection{Архитектурные подходы и обоснование выбранного стека}

В проектировании прикладных информационных систем можно выделить два полярных подхода: монолитный сценарий с жесткой связностью компонентов и модульный сценарий с разделением ответственности. Для учебного проекта второго курса/третьего курса (в зависимости от учебного плана) предпочтителен второй вариант, поскольку он лучше демонстрирует инженерную зрелость и позволяет управлять сложностью разработки \cite{sommerville,designPatterns}.

С учетом поставленных задач в \textit{Gamified Task CRM} используется кроссплатформенный UI-фреймворк Flutter, язык Dart, механизм управления состоянием Riverpod, декларативная маршрутизация GoRouter и backend-платформа Supabase. Ниже приведено сравнительное обоснование такого выбора.

\begin{table}[h!]
\caption{Сравнение ключевых технологических решений проекта}
\centering
\begin{tabular}{|p{4cm}|p{4.5cm}|p{5.7cm}|}
\hline
\textbf{Компонент} & \textbf{Рассмотренные варианты} & \textbf{Причина выбора в проекте} \\
\hline
UI-фреймворк & Flutter, React + TypeScript, Vue & Единая кодовая база, быстрый UI-прототипинг, зрелая экосистема для web и мобильного расширения \cite{flutterDocs} \\
\hline
Управление состоянием & setState, BLoC, Riverpod & Декларативность, тестируемость, удобная композиция провайдеров, меньшая связность \cite{riverpodDocs} \\
\hline
Маршрутизация & Navigator 1.0, auto\_route, GoRouter & Прозрачные redirect-правила, удобная интеграция с авторизацией и deep-link сценариями \cite{goRouterDocs} \\
\hline
Backend & Firebase, Supabase, custom API & PostgreSQL-модель данных, RLS-политики, SQL-миграции и удобный локальный/облачный контур \cite{supabaseDocs,postgresDocs} \\
\hline
\end{tabular}
\label{tab:tech-comparison}
\end{table}

Данные таблицы \ref{tab:tech-comparison} показывают, что выбор стека основан не на популярности инструментов, а на соответствии требованиям проекта: управляемость состояния, прозрачная маршрутизация, поддержка реляционной модели и доступность практик безопасности. Для курсовой работы такой подход является корректным, поскольку демонстрирует умение сопоставлять техническое решение с предметной задачей.

С архитектурной точки зрения применим принцип \textit{feature-first organization}, когда код группируется по функциональным модулям (dashboard, tasks, crm, reports, profile, ai), а не по техническим слоям в изоляции. Это решение улучшает локальность изменений: добавление нового поведения в одном модуле требует минимального вмешательства в соседние. При этом общие инфраструктурные элементы (конфигурация, база, авторизация) остаются централизованными и переиспользуемыми \cite{repoReadme,patternsBook}.

Для формального описания эксплуатационной результативности можно использовать показатель покрытия ключевых сценариев:
\begin{equation}
Q = \frac{N_{pass}}{N_{total}}.
\label{eq:quality}
\end{equation}

где $N_{pass}$ --- количество успешно пройденных проверок, а $N_{total}$ --- общее количество запущенных проверок в контрольном контуре. Показатель \eqref{eq:quality} не исчерпывает понятие качества, но служит удобной метрикой для оперативного мониторинга технической стабильности при развитии проекта \cite{testingBook}.

\subsection{Теоретические основы проектирования данных и клиентских сценариев}

Проектирование модели данных является центральной задачей для любой CRM-системы, поскольку именно данные связывают пользовательские действия, бизнес-логику и аналитические отчеты. В рамках рассматриваемого проекта ключевая логическая связь задается через отношение \textit{tasks.client\_id $\rightarrow$ crm\_clients.id}. Такая связь допускает как привязанные, так и автономные задачи, что отражает реальный процесс работы: часть задач связана с конкретным клиентом, часть относится к внутренним активностям команды \cite{repoReadme,migrationClientId}.

При построении реляционной модели для учебного проекта целесообразно придерживаться следующих правил:
\begin{enumerate}
  \item Каждая сущность должна иметь устойчивый первичный идентификатор.
  \item Внешние связи должны быть явными и проверяемыми на уровне схемы данных.
  \item Поля с неопределенным значением должны иметь корректную nullable-семантику.
  \item Дата-время в операционных таблицах должна сопровождаться полями создания и обновления.
\end{enumerate}

Указанные правила соответствуют базовым принципам проектирования информационных систем и существенно упрощают последующую интеграцию прикладной логики с хранилищем. Например, наличие полей \textit{created\_at} и \textit{updated\_at} позволяет организовать сортировку, аудит изменений и формирование периодических отчетов без дополнительной денормализации \cite{dbBook2021,postgresDocs}.

С точки зрения пользовательских сценариев важна не только структурная корректность данных, но и целостность переходов между экранами. Типовой процесс CRM-пользователя можно представить в виде цепочки: вход в систему, создание или выбор клиента, постановка задачи, контроль исполнения, просмотр сводной информации. Наличие разрывов в этой цепочке снижает практическую ценность приложения, поскольку пользователь вынужден выполнять внешние действия вне системы.

В теории UX-проектирования это описывается как требование \textit{continuity of workflow}: пользователь должен завершить ключевой сценарий без неоправданных переключений контекста \cite{uxBook2023}. В проекте \textit{Gamified Task CRM} данное требование поддерживается за счет единой оболочки с нижней навигацией и согласованных форм добавления сущностей. С инженерной точки зрения это также уменьшает вероятность ошибок, связанных с дублированием логики в независимых экранах.

На рисунке \ref{fig:architecture} показана обобщенная схема модульного взаимодействия в системе.

\begin{figure}[h!]
\centering
\fbox{%
\begin{minipage}{0.9\textwidth}
\centering
UI-модули (Dashboard / Tasks / CRM / Reports / Profile / AI) \\
$\downarrow$ \\
Провайдеры состояния (Riverpod Notifiers) \\
$\downarrow$ \\
Сервисы данных (Supabase API, Auth, локальный demo-контур) \\
$\downarrow$ \\
PostgreSQL-таблицы и политики доступа
\end{minipage}}
\caption{Обобщенная архитектурная схема проекта}
\label{fig:architecture}
\end{figure}

Как показано на рисунке \ref{fig:architecture}, модульность достигается не только разбиением интерфейса, но и четким разделением уровней ответственности: UI, состояние, сервисы, хранилище. Такой подход упрощает сопровождение и снижает риск каскадных дефектов при изменении одного из компонентов.

\subsection{Качество программного продукта: тестирование, надежность, безопасность}

Качество программного продукта в инженерной практике интерпретируется как совокупность нескольких свойств: функциональная корректность, устойчивость к ошибкам, эксплуатационная предсказуемость и безопасность. Для курсового проекта важно показать, что эти свойства учитываются не декларативно, а через конкретные процедуры проверки \cite{testingBook,owaspM10}.

В рассматриваемой системе используется многоуровневый подход к верификации:
\begin{enumerate}
  \item \textbf{Статический анализ} --- выявление потенциальных ошибок и нарушений стиля до выполнения приложения.
  \item \textbf{Unit/Widget-тесты} --- проверка отдельных компонентов и UI-поведения в изолированной среде.
  \item \textbf{Integration smoke} --- контроль целостности основного бизнес-потока после изменений.
  \item \textbf{Web E2E} --- сквозная проверка критичного сценария в среде, максимально близкой к эксплуатационной.
\end{enumerate}

Теоретически такой набор тестов соответствует принципу \textit{test pyramid}: широкая база быстрых проверок на нижних уровнях и ограниченное число более дорогих сквозных сценариев на верхнем уровне \cite{testingBook}. Для учебного проекта это особенно важно, так как позволяет соблюдать баланс между глубиной проверки и временем выполнения.

Отдельную роль играет обработка исключительных ситуаций. Даже корректно спроектированная архитектура теряет практическую ценность, если пользователь получает неинформативные сообщения об ошибках. В рамках проекта используется нормализация технических backend-ошибок в человеко-понятные сообщения (например, отсутствие таблицы, проблема RLS, истекшая сессия). Теоретически это относится к паттерну \textit{error translation boundary}, когда низкоуровневые ошибки преобразуются на границе слоя сервиса в формат, пригодный для пользовательского интерфейса \cite{cleanArchitecture,patternsBook}.

С позиции информационной безопасности в курсовом проекте применяются базовые, но принципиальные практики:
\begin{itemize}
  \item вынесение секретов из кода в параметры окружения;
  \item разделение demo и production режимов;
  \item ограничение доступа к данным через политики уровня строк;
  \item контроль аутентификации перед операциями записи.
\end{itemize}

Эти меры соответствуют рекомендациям OWASP по предотвращению типичных ошибок конфигурации и контроля доступа в web-приложениях \cite{owaspM4,owaspM1}. Для учебной разработки их достаточно, чтобы продемонстрировать осознанный подход к защищенной реализации.

\subsection{Теоретические выводы по разделу}

Проведенный анализ показывает, что разработка CRM-ориентированного приложения требует системного сочетания трех групп решений: предметных (что делает система), архитектурных (как организован код) и эксплуатационных (как подтверждается корректность и устойчивость). Игнорирование любой из этих групп приводит либо к функционально неполному продукту, либо к трудно сопровождаемой реализации.

Для темы курсового проекта обоснован выбор Flutter-стека с Riverpod, GoRouter и Supabase, поскольку он обеспечивает:
\begin{itemize}
  \item модульную организацию кода и масштабируемость;
  \item прозрачные пользовательские сценарии;
  \item реляционную модель данных с контролируемыми связями;
  \item возможность поэтапной верификации качества.
\end{itemize}

Теоретическая часть также подтверждает, что качество итогового решения определяется не только полнотой функционала, но и зрелостью инженерных практик: корректной конфигурацией окружения, продуманной обработкой ошибок, повторяемыми тестовыми сценариями и соблюдением минимальных требований безопасности.

Сформированные в разделе положения являются методической базой для практической главы, где рассматривается реализация \textit{Gamified Task CRM}, описание архитектуры в коде, структура данных, сценарии использования и результаты проверки работоспособности.

\clearpage
\section{ПРАКТИЧЕСКАЯ РЕАЛИЗАЦИЯ ПРОЕКТА GAMIFIED TASK CRM}

\subsection{Общая структура реализации и конфигурация среды}

Практическая реализация проекта выполнена в репозитории \textit{gamified\_app} и ориентирована на web-режим запуска Flutter-приложения. Входной точкой системы является файл \texttt{lib/main.dart}, в котором организованы базовые этапы инициализации: подготовка Flutter binding, инициализация локального хранилища Hive, регистрация адаптеров моделей, запуск сервиса аутентификации и подъем корневого виджета приложения в \texttt{ProviderScope}. Такое построение старта соответствует требованию предсказуемого bootstrap-процесса и исключает скрытые неявные зависимости между модулями \cite{repoReadme}.

С точки зрения конфигурации окружения проект использует параметры \texttt{--dart-define}. К обязательным для backend-first контура относятся \texttt{SUPABASE\_URL} и \texttt{SUPABASE\_ANON\_KEY}; для AI-функций подключается \texttt{MINIMAX\_API\_KEY}. Дополнительно предусмотрены флаги режимов: \texttt{DEMO\_MODE=true} и \texttt{AI\_DEMO\_MODE=true}. Эта схема важна как с эксплуатационной, так и с безопасностной точки зрения: секреты не хранятся в исходном коде, а режимы исполнения разделяются явно \cite{repoReadme,owaspM4}.

Маршрутизация организована через \texttt{GoRouter}. В приложении заданы два главных маршрута: \texttt{/login} и \texttt{/app}. Логика \texttt{redirect} обеспечивает автоматическое перенаправление неаутентифицированного пользователя на экран входа и блокирует доступ в закрытую часть интерфейса без валидной сессии. При этом для авторизованного пользователя обратный переход на \texttt{/login} также блокируется. Такой механизм устраняет класс ошибок, когда часть защищенного функционала случайно остается доступной по прямому URL \cite{goRouterDocs,owaspM1}.

Пользовательская оболочка реализована в \texttt{RootLayout} как \texttt{IndexedStack} с пятью вкладками: главная панель, задачи, CRM, отчеты и профиль. Индексная модель хранения экранов позволяет сохранять локальное состояние вкладок и ускоряет переключение между ними, что особенно заметно при частой навигации между задачами и клиентскими карточками. Дополнительные действия создания сущностей (задача, клиент) вынесены в контекстные формы через \texttt{FloatingActionButton} и \texttt{ModalBottomSheet}, что делает поток работы компактным и однотипным во всех основных модулях.

В практической части курсового проекта важно подчеркнуть, что архитектурные решения выбирались исходя из учебной цели: показать контролируемую, расширяемую систему, а не точечный прототип. Поэтому структура проекта поддерживает функциональную декомпозицию по доменным блокам:
\begin{itemize}
  \item \texttt{lib/features/dashboard/presentation/};
  \item \texttt{lib/features/tasks/presentation/};
  \item \texttt{lib/features/crm/presentation/};
  \item \texttt{lib/features/reports/presentation/};
  \item \texttt{lib/features/profile/presentation/};
  \item \texttt{lib/features/ai/presentation/}.
\end{itemize}

Такое деление улучшает поддерживаемость: разработчик работает в рамках конкретной предметной зоны, не затрагивая нерелевантные части системы. Кроме того, модульная структура облегчает контроль требований при защите курсового проекта, поскольку каждую функциональную область можно объяснить отдельно: назначение, входные данные, основной сценарий, ограничения и тестовый контур.

\subsection{Модель данных и backend-first интеграция}

Практическая реализация данных строится на двух взаимосвязанных источниках: локальное хранилище (demo-режим) и облачный backend (рабочий режим). На уровне клиентского кода модель задачи описывается классом \texttt{Task}, включающим поля \texttt{id}, \texttt{title}, \texttt{description}, \texttt{priority}, \texttt{category}, \texttt{clientId}, \texttt{dueDate}, \texttt{createdAt}, \texttt{isCompleted}. Для CRM используется модель \texttt{Client} с идентификатором, контактными атрибутами, статусом, источником и датой создания.

Для backend-контура в README проекта зафиксирован минимальный состав таблиц: \texttt{public.tasks}, \texttt{public.crm\_clients}, \texttt{public.user\_stats}. Ключевой миграцией выступает добавление связи \texttt{tasks.client\_id} с таблицей клиентов, что обеспечивает трассируемость сценария \textit{задача, привязанная к конкретному клиенту} \cite{repoReadme,migrationClientId,bootstrapSchema}.

С инженерной точки зрения важна не только структура таблиц, но и дисциплина применения схемы. В проекте предусмотрен идемпотентный SQL-скрипт bootstrap, создающий/обновляющий базовые таблицы, внешние ключи, индексы и триггеры \texttt{updated\_at}. Такой подход снижает риск дрейфа схемы между окружениями и упрощает повторную настройку стенда перед защитой.

В таблице \ref{tab:data-model} приведена укрупненная характеристика ключевых сущностей системы.

\begin{table}[h!]
\caption{Ключевые сущности и их роль в системе}
\centering
\begin{tabular}{|p{3.5cm}|p{5.2cm}|p{5.5cm}|}
\hline
\textbf{Сущность} & \textbf{Основные поля} & \textbf{Назначение в бизнес-потоке} \\
\hline
Task & id, title, priority, status, due\_date, client\_id & Планирование и контроль операционных действий пользователя, связь с клиентом при необходимости \\
\hline
CRM Client & id, first\_name, last\_name, email, phone, status, source & Хранение клиентской базы и контекста взаимодействия, сегментация контактов \\
\hline
User Stats & xp, level, streak, last\_task\_date & Геймифицированные метрики активности и мотивации пользователя \\
\hline
\end{tabular}
\label{tab:data-model}
\end{table}

Как видно из таблицы \ref{tab:data-model}, модель данных не перегружена второстепенными сущностями. Это осознанное решение, направленное на устойчивость базового сценария. Слишком раннее усложнение схемы (например, введение десятков справочников) в учебном проекте обычно приводит к росту технического долга и снижению качества тестирования.

Важной частью реализации является обработка ошибок backend-взаимодействия. В \texttt{main.dart} реализован метод нормализации исключений в пользовательские сообщения. Примеры покрываемых случаев: отсутствие таблиц в Supabase, ошибки политики RLS, некорректный UUID, истекший JWT-токен. Такая прослойка уменьшает когнитивную нагрузку на пользователя и ускоряет диагностику проблем во время демонстрации проекта.

Если формализовать надежность клиент-серверного обмена на уровне сессии, можно использовать показатель доступности операций:
\begin{equation}
A = \frac{N_{ok}}{N_{ok} + N_{err}},
\label{eq:availability}
\end{equation}

где $N_{ok}$ --- количество успешно завершенных операций, $N_{err}$ --- количество операций, завершившихся ошибкой. В практической разработке показатель \eqref{eq:availability} используется как оперативная оценка стабильности интеграции после изменений в схеме или политике доступа.

\subsection{Реализация пользовательских сценариев и модульной логики}

Ключевой прикладной сценарий проекта включает последовательность действий: вход пользователя в систему, создание карточки клиента, постановка задачи с привязкой к клиенту, проверка отображения результата в списке задач и в деталях клиента. Этот сценарий отражает основной смысл CRM-подхода: управление работой строится не изолированно, а в контексте конкретных контактов и сделок \cite{repoReadme,docsDefenseScript}.

В модуле задач реализованы базовые CRUD-операции, приоритеты, категории и сроки выполнения. Привязка к клиенту сделана nullable-полем, что поддерживает две бизнес-модели: задачи без клиентского контекста (внутренние операции) и задачи с клиентским контекстом (внешние взаимодействия). В CRM-модуле реализовано управление карточками клиентов и их статусами, что позволяет отслеживать жизненный цикл взаимодействия от лида до активного клиента.

Модуль отчетности в текущем контуре предназначен для представления агрегированных показателей, а не для сложной BI-аналитики. Такой уровень детализации соответствует масштабу курсового проекта: пользователь получает базовую оценку результативности, при этом система остается легкой для сопровождения. При дальнейшем развитии проекта модуль отчетов можно расширить когортным анализом, сегментацией по источникам лидов и динамикой конверсии по периодам.

Важное практическое свойство реализации --- единый UX-паттерн ввода данных через нижние модальные формы. Он снижает вероятность сценарных расхождений между модулями и ускоряет обучение пользователя работе с системой. В учебном контексте это также упрощает проведение защиты: демонстрация строится на повторяемых действиях, а не на случайных переходах по сложному интерфейсу.

В таблице \ref{tab:scenarios} приведены основные сценарии и ожидаемые результаты.

\begin{table}[h!]
\caption{Ключевые пользовательские сценарии и ожидаемый эффект}
\centering
\begin{tabular}{|p{4.5cm}|p{5.2cm}|p{4.5cm}|}
\hline
\textbf{Сценарий} & \textbf{Проверяемое поведение} & \textbf{Практический результат} \\
\hline
Авторизация пользователя & Redirect между \texttt{/login} и \texttt{/app} в зависимости от сессии & Защита закрытого контура и корректный старт работы \\
\hline
Создание клиента & Валидация полей и сохранение карточки в локальный/backend контур & Актуальная CRM-база для планирования задач \\
\hline
Создание задачи с client\_id & Сохранение связи \textit{задача-клиент} и отображение в двух представлениях & Трассируемость операционной активности \\
\hline
Просмотр отчета & Агрегация базовых показателей по активности & Быстрая обратная связь по прогрессу \\
\hline
Обработка ошибок backend & Нормализация технической ошибки в понятное сообщение & Снижение времени диагностики и доработки \\
\hline
\end{tabular}
\label{tab:scenarios}
\end{table}

Как видно из таблицы \ref{tab:scenarios}, практическая реализация охватывает полный минимальный контур пользовательской ценности: вход в систему, данные клиентов, постановка задач и контроль результата.

Отдельно следует отметить геймифицированную составляющую проекта. В модели \texttt{UserStats} учитываются показатели XP, уровень и streak. Несмотря на упрощенный характер этой подсистемы, она решает важную задачу мотивации пользователя к регулярному выполнению задач. С точки зрения продуктовой логики это оправдано: даже простые механики прогресса повышают вовлеченность в системах персональной продуктивности \cite{productivityBook}.

\subsection{Верификация качества и результаты практической проверки}

Проверка качества в практической главе организована как последовательность обязательных шагов. На первом уровне выполняется \texttt{flutter analyze} для выявления статических проблем. На втором уровне запускаются автоматизированные тесты: \texttt{flutter test} и \texttt{flutter test --coverage}. На третьем уровне применяются скрипты контроля покрытия и интеграционного smoke-потока. На четвертом уровне предусмотрен web E2E-сценарий с реальной аутентификацией и backend-операциями \cite{repoReadme}.

С методической точки зрения такая организация проверки дает несколько преимуществ:
\begin{itemize}
  \item раннее обнаружение синтаксических и типовых ошибок;
  \item контроль корректности доменной логики на уровне модулей;
  \item подтверждение целостности основного пользовательского сценария;
  \item снижение риска регрессий перед защитой проекта.
\end{itemize}

В таблице \ref{tab:quality-checks} представлена сводка контрольного контура качества.

\begin{table}[h!]
\caption{Контур проверок качества проекта}
\centering
\begin{tabular}{|p{4.2cm}|p{5.5cm}|p{4.5cm}|}
\hline
\textbf{Инструмент} & \textbf{Что проверяется} & \textbf{Назначение} \\
\hline
\texttt{flutter analyze} & Предупреждения, потенциальные ошибки типов и API-вызовов & Статическая устойчивость кода \\
\hline
\texttt{flutter test} & Unit и widget сценарии & Корректность логики и UI-компонентов \\
\hline
\texttt{flutter test --coverage} + gate скрипты & Полнота тестового охвата runtime-критичного кода & Контроль инженерной дисциплины \\
\hline
Integration smoke script & Базовый пользовательский поток в интеграционной среде & Быстрая проверка работоспособности после изменений \\
\hline
Web Supabase E2E script & Сквозной сценарий с реальным backend и авторизацией & Подтверждение эксплуатационной готовности \\
\hline
\end{tabular}
\label{tab:quality-checks}
\end{table}

Результаты практической проверки показывают, что выбранный набор инструментов является достаточным для учебного проекта, где важны не абсолютные метрики enterprise-уровня, а воспроизводимость и аргументированность качества. Наличие четко описанных команд запуска также облегчает приемку преподавателем: сценарии можно повторить без ручного вмешательства в код.

Для формализации проверки к отчету добавлены приложения: полный набор команд запуска и проверки (Приложение 1), структура ключевых директорий (Приложение 2), а также таблица приемочных критериев (Приложение 3). Ссылки на эти материалы позволяют быстро сопоставить теоретическое описание и практический контур реализации.

При этом необходимо зафиксировать ограничения. Во-первых, в проекте пока не реализован полный контур нагрузочного тестирования, поэтому поведение при высоком числе одновременных пользователей оценивается косвенно. Во-вторых, модуль отчетности сфокусирован на базовых сводках и не покрывает сложные BI-сценарии. В-третьих, AI-подсистема рассматривается как вспомогательный модуль и не является критичным контуром бизнес-логики.

Указанные ограничения не снижают учебную ценность проекта, поскольку главная цель курсовой работы состоит в демонстрации корректно спроектированного и реализованного базового решения. Более того, наличие явно сформулированных ограничений является признаком зрелого инженерного подхода, когда команда осознает текущие границы системы и планирует дальнейшее развитие.

\subsection{Экономия трудозатрат, риски и направления развития}

Практический эффект от реализации \textit{Gamified Task CRM} оценивается через сокращение ручных операций и упорядочивание взаимодействия с клиентами. В простейшем сценарии пользователь получает единое рабочее пространство вместо разрозненных таблиц, мессенджеров и заметок. Это снижает потери контекста и улучшает управляемость ежедневных задач.

С точки зрения проектных рисков можно выделить три основные группы:
\begin{enumerate}
  \item \textbf{Конфигурационные риски}: ошибки в переменных окружения, несовпадение схемы БД, некорректные политики доступа.
  \item \textbf{Продуктовые риски}: избыточная сложность интерфейса при масштабировании функционала.
  \item \textbf{Технические риски}: накопление технического долга при быстром добавлении новых модулей без рефакторинга.
\end{enumerate}

Для управления этими рисками в проекте применяются базовые меры: bootstrap-скрипт схемы, автоматизированные проверки, модульная структура и явные runtime-флаги режимов. В дальнейшем целесообразно усилить практику за счет централизованной телеметрии, мониторинга ошибок и формализации release-процесса.

Приоритетные направления развития после завершения курсового проекта могут быть сформулированы следующим образом:
\begin{itemize}
  \item расширение аналитического блока (воронка, сегментация, динамика конверсии);
  \item добавление ролевой модели пользователей и разграничения прав;
  \item внедрение уведомлений по срокам и событиям;
  \item развитие мобильного канала на базе той же кодовой базы Flutter;
  \item усиление quality-контура за счет дополнительных интеграционных и контрактных тестов.
\end{itemize}

Суммарно практическая реализация подтверждает достижение промежуточного целевого состояния: создан работоспособный, структурно устойчивый и расширяемый программный продукт, соответствующий учебным требованиям по предметной области, архитектуре, данным и качеству выполнения.

\clearpage
\section*{ЗАКЛЮЧЕНИЕ}
\addcontentsline{toc}{section}{ЗАКЛЮЧЕНИЕ}

В ходе выполнения курсового проекта была достигнута поставленная цель --- разработано и обосновано программное решение \textit{Gamified Task CRM (Flutter Web)}, ориентированное на управление задачами и клиентскими взаимодействиями в едином цифровом контуре. Работа показала, что даже в рамках учебного проекта возможно реализовать систему, в которой объединены предметная логика CRM, структурированная архитектура, backend-first подход к данным и формализованный контроль качества.

Первая задача, связанная с анализом теоретических и технологических основ разработки CRM-ориентированного приложения, выполнена в полном объеме. В теоретической части рассмотрены назначение CRM-систем, требования к функциональному и нефункциональному контуру, подходы к архитектуре, принципы проектирования данных и базовые практики надежности и безопасности. На этой основе обоснован выбор стека Flutter + Riverpod + GoRouter + Supabase как рационального решения для проекта с требованиями к модульности, прозрачной маршрутизации и реляционной модели данных.

Вторая задача, предусматривающая проектирование логической структуры решения, также выполнена. В проекте выделены функциональные модули интерфейса (dashboard, tasks, crm, reports, profile, ai), определены границы ответственности, реализована единая оболочка навигации и согласованный шаблон пользовательских действий. На уровне данных зафиксирована ключевая связь \textit{задача-клиент} через \texttt{client\_id}, что обеспечивает трассируемость операционной активности. Подобная структура позволяет последовательно масштабировать систему без радикального пересмотра архитектуры.

Третья задача, связанная с практической реализацией, решена путем построения рабочего web-приложения на Flutter с поддержкой авторизации, управления клиентами и задачами, базовой аналитики и режимов исполнения. Реализация содержит как backend-first контур для реального использования, так и demo-режим для воспроизводимой учебной демонстрации. Такой подход позволил совместить практичность и устойчивость: система работает с внешним backend, но при необходимости может быть показана в автономном режиме без потери основной функциональности.

Четвертая задача --- проверка корректности и устойчивости решения --- выполнена через многоуровневый quality-контур: статический анализ, автоматизированные тесты, smoke-проверки и web E2E-сценарий. В результате подтверждена работоспособность ключевого пользовательского процесса \textit{Login $\rightarrow$ Create Client $\rightarrow$ Create Task $\rightarrow$ Verify}. Применение формализованных проверок повысило доверие к результатам разработки и обеспечило воспроизводимость демонстрации перед защитой.

Итоговые результаты курсового проекта можно свести к следующим положениям:
\begin{enumerate}
  \item Создано прикладное CRM-ориентированное приложение с модульной архитектурой и четкой декомпозицией функционала.
  \item Обеспечена связность операционных сущностей, включая корректную модель данных для задач, клиентов и пользовательской статистики.
  \item Реализован безопасный конфигурационный подход через переменные окружения и разграничение demo/production режимов.
  \item Сформирован практический контур верификации качества, достаточный для учебной и предэксплуатационной проверки.
  \item Подготовлена база для дальнейшего развития проекта без нарушения базовых архитектурных принципов.
\end{enumerate}

Практическая значимость работы заключается в том, что разработанное решение может использоваться как учебный стенд полного цикла разработки программного обеспечения. Проект позволяет демонстрировать не только интерфейсную часть, но и инженерные дисциплины: организацию данных, контроль качества, управление конфигурацией, обработку ошибок, подготовку к деплою. Это делает работу применимой как для образовательных задач, так и для стартовых прикладных сценариев малого бизнеса или учебных лабораторий.

Одновременно следует обозначить ограничения текущей версии. В системе пока отсутствуют расширенные аналитические инструменты уровня BI, полноценная ролевая модель с многоуровневым разграничением прав и глубокий нагрузочный контур. Кроме того, AI-модуль играет вспомогательную роль и требует дополнительного исследования в части качества ответов и прикладной полезности для CRM-процессов. Тем не менее эти ограничения являются ожидаемыми для масштаба курсового проекта и не противоречат достижению его основной цели.

С учетом полученных результатов можно предложить следующие направления дальнейшего развития:
\begin{enumerate}
  \item Расширить отчетный блок: добавить воронку продаж, периодическую динамику и сегментацию клиентской базы.
  \item Ввести ролевую модель доступа и аудит критичных операций.
  \item Реализовать сервис уведомлений по дедлайнам задач и событиям CRM.
  \item Усилить качество через дополнительные интеграционные и контрактные тесты.
  \item Масштабировать продукт в сторону mobile-first сценариев на общей Flutter-кодовой базе.
\end{enumerate}

Подводя итог, можно заключить, что курсовой проект демонстрирует достижение поставленной цели и решение всех сформулированных задач. Разработанное программное обеспечение обладает практической применимостью, соответствует требованиям к структуре и содержанию учебной работы и служит устойчивой основой для последующего инженерного развития.


\clearpage
\clearpage
\section*{СПИСОК ИСПОЛЬЗОВАННОЙ ЛИТЕРАТУРЫ (APA)}
\addcontentsline{toc}{section}{СПИСОК ИСПОЛЬЗОВАННОЙ ЛИТЕРАТУРЫ (APA)}

\begingroup
\setlength{\bibhang}{1.25cm}
\setlength{\bibsep}{0.4\baselineskip}
\begin{thebibliography}{99}

\bibitem[Antonova \& Smirnov(2024)]{metrics2024}
Antonova, I. I., \& Smirnov, V. A. (2024). \textit{Statisticheskie metody v upravlenii kachestvom}. Yurait.

\bibitem[Buttle \& Maklan(2022)]{crmBook}
Buttle, F., \& Maklan, S. (2022). \textit{Customer relationship management: Concepts and technologies} (4th ed.). Routledge.

\bibitem[Chistov et~al.(2025)]{isDesign2025}
Chistov, D. V., Melnikov, P. P., Zolotaryuk, A. V., \& Nicheporuk, N. B. (2025). \textit{Proektirovanie informacionnyh sistem: Uchebnik dlya SPO}. Yurait.

\bibitem[Dart Team(2026)]{dartDocs}
Dart Team. (2026). \textit{Dart documentation}. \url{https://dart.dev/guides}

\bibitem[Date(2022)]{dataDesign2022}
Date, C. J. (2022). \textit{Database design and relational theory} (2nd ed.). O'Reilly Media.

\bibitem[Eyal(2021)]{productivityBook}
Eyal, N. (2021). \textit{Hooked: How to build habit-forming products}. Portfolio.

\bibitem[Freeman \& Robson(2021)]{designPatterns}
Freeman, E., \& Robson, E. (2021). \textit{Head first design patterns} (2nd ed.). O'Reilly Media.

\bibitem[Flutter Team(2026a)]{flutterDocs}
Flutter Team. (2026a). \textit{Flutter documentation}. \url{https://docs.flutter.dev}

\bibitem[Flutter Team(2026b)]{goRouterDocs}
Flutter Team. (2026b). \textit{GoRouter package documentation}. \url{https://pub.dev/packages/go_router}

\bibitem[Humble \& Farley(2022)]{softwareEconomics2022}
Humble, J., \& Farley, D. (2022). \textit{Continuous delivery and software economics in practice}. Addison-Wesley.

\bibitem[Kerzner(2022)]{projectManagementBook}
Kerzner, H. (2022). \textit{Project management: A systems approach to planning, scheduling, and controlling}. Wiley.

\bibitem[Krug(2023)]{uxBook2023}
Krug, S. (2023). \textit{Don't make me think, revisited for modern products}. New Riders.

\bibitem[Martin(2021)]{cleanArchitecture}
Martin, R. C. (2021). \textit{Clean architecture: A craftsman's guide to software structure and design}. Prentice Hall.

\bibitem[OWASP Foundation(2021a)]{owaspM1}
OWASP Foundation. (2021a). \textit{OWASP Top 10 2021: Broken access control}. \url{https://owasp.org/Top10/A01_2021-Broken_Access_Control/}

\bibitem[OWASP Foundation(2021b)]{owaspM4}
OWASP Foundation. (2021b). \textit{OWASP Top 10 2021: Insecure design}. \url{https://owasp.org/Top10/A04_2021-Insecure_Design/}

\bibitem[OWASP Foundation(2021c)]{owaspM10}
OWASP Foundation. (2021c). \textit{OWASP Top 10 2021}. \url{https://owasp.org/Top10/}

\bibitem[Pecorelli et~al.(2023)]{testingBook}
Pecorelli, A., et al. (2023). \textit{Testing Flutter applications}. Manning.

\bibitem[PostgreSQL Global Development Group(2025)]{postgresDocs}
PostgreSQL Global Development Group. (2025). \textit{PostgreSQL 16 documentation}. \url{https://www.postgresql.org/docs/16/}

\bibitem[Project Team(2026a)]{repoReadme}
Project Team. (2026a). \textit{Gamified Task CRM (Flutter Web): README and runtime guide}. Internal project documentation.

\bibitem[Project Team(2026b)]{docsIndex}
Project Team. (2026b). \textit{Docs index for current runtime}. Internal project documentation.

\bibitem[Project Team(2026c)]{docsDefenseScript}
Project Team. (2026c). \textit{Defense demo script}. Internal project documentation.

\bibitem[Project Team(2026d)]{migrationClientId}
Project Team. (2026d). \textit{Migration: Add client\_id to public.tasks}. SQL migration document.

\bibitem[Project Team(2026e)]{bootstrapSchema}
Project Team. (2026e). \textit{Bootstrap runtime schema for Supabase}. SQL migration document.

\bibitem[Richards \& Ford(2022)]{patternsBook}
Richards, M., \& Ford, N. (2022). \textit{Fundamentals of software architecture}. O'Reilly Media.

\bibitem[Rousselet(2026)]{riverpodDocs}
Rousselet, R. (2026). \textit{Riverpod documentation}. \url{https://riverpod.dev}

\bibitem[Satzinger et~al.(2023)]{isDesign2023}
Satzinger, J., Jackson, R., \& Burd, S. (2023). \textit{Systems analysis and design in a changing world} (8th ed.). Cengage.

\bibitem[Silberschatz et~al.(2021)]{dbBook2021}
Silberschatz, A., Korth, H., \& Sudarshan, S. (2021). \textit{Database system concepts} (7th ed.). McGraw-Hill.

\bibitem[Sommerville(2022)]{sommerville}
Sommerville, I. (2022). \textit{Software engineering} (11th ed.). Pearson.

\bibitem[Supabase Inc.(2026)]{supabaseDocs}
Supabase Inc. (2026). \textit{Supabase documentation}. \url{https://supabase.com/docs}

\end{thebibliography}
\endgroup


\clearpage
\pagestyle{empty}
\appendixsection{1}{Команды запуска, проверки и деплоя}

\textbf{Базовый запуск проекта (web):}
\begin{verbatim}
flutter pub get
flutter run -d chrome
\end{verbatim}

\textbf{Запуск backend-first с параметрами окружения:}
\begin{verbatim}
flutter run -d chrome \
  --dart-define=DEMO_MODE=false \
  --dart-define=SUPABASE_URL=... \
  --dart-define=SUPABASE_ANON_KEY=... \
  --dart-define=MINIMAX_API_KEY=...
\end{verbatim}

\textbf{Проверки качества:}
\begin{verbatim}
flutter analyze
flutter test
flutter test --coverage
bash scripts/ci/report_full_lib_coverage.sh coverage/lcov.info
bash scripts/ci/check_coverage_gate.sh coverage/lcov.info 55 \
  scripts/ci/runtime_coverage_scope.txt
bash scripts/ci/run_integration_smoke.sh
\end{verbatim}

\textbf{Web E2E с Supabase:}
\begin{verbatim}
SUPABASE_URL=... \
SUPABASE_ANON_KEY=... \
E2E_EMAIL=... \
E2E_PASSWORD=... \
E2E_DEVICE=chrome \
bash scripts/ci/run_web_supabase_e2e.sh
\end{verbatim}

\appendixsection{2}{Структура ключевых директорий проекта}

\begin{verbatim}
lib/
  main.dart
  app/
    root_layout.dart
  features/
    dashboard/presentation/
    tasks/presentation/
    crm/presentation/
    reports/presentation/
    profile/presentation/
    ai/presentation/
  providers.dart
  providers_crm.dart

docs/
  README.md
  migrations/
    2026-03-04-bootstrap-runtime-schema.sql
    2026-03-04-add-client-id-to-public-tasks.sql
  reports/
    2026-03-04-defense-demo-script.md
\end{verbatim}

\appendixsection{3}{Фрагменты критериев приемки функционала}

\begin{longtable}{|p{5.5cm}|p{7.5cm}|}
\hline
\textbf{Критерий} & \textbf{Проверка} \\
\hline
Авторизация и redirect & Неавторизованный пользователь не попадает в \texttt{/app}; после входа открывается рабочий контур \\
\hline
CRUD клиентов & Создание, изменение и удаление карточек выполняется без ошибок в выбранном режиме \\
\hline
CRUD задач & Добавление и обновление задач отражается в списке и сохраняется в хранилище \\
\hline
Связь task-client & При указании \texttt{client\_id} задача отображается в карточке соответствующего клиента \\
\hline
Базовые отчеты & Модуль отчетности показывает агрегированные показатели по активности \\
\hline
Стабильность сборки & \texttt{flutter analyze} и базовые тесты выполняются без блокирующих ошибок \\
\hline
\end{longtable}


\end{document}
